
\documentclass{article}
\usepackage{colortbl}
\usepackage{makecell}
\usepackage{multirow}
\usepackage{supertabular}

\begin{document}

\newcounter{utterance}

\centering \large Interaction Transcript for game `dond', experiment `coop', episode 4 with deepseek{-}v4{-}pro.
\vspace{24pt}

{ \footnotesize  \setcounter{utterance}{1}
\setlength{\tabcolsep}{0pt}
\begin{supertabular}{c@{$\;$}|p{.15\linewidth}@{}p{.15\linewidth}p{.15\linewidth}p{.15\linewidth}p{.15\linewidth}p{.15\linewidth}}
    \# & $\;$A & \multicolumn{4}{c}{Game Master} & $\;\:$B\\
    \hline

    \theutterance \stepcounter{utterance}  
    & & \multicolumn{4}{p{0.6\linewidth}}{
        \cellcolor[rgb]{0.9,0.9,0.9}{
            \makecell[{{p{\linewidth}}}]{
                \texttt{\tiny{[P1$\langle$GM]}}
                \texttt{Vous jouez à un jeu de négociation coopérative dans lequel vous devez vous entendre avec un autre joueur sur la façon de diviser un ensemble d'objets entre vous.} \\
\\ 
\texttt{Règles:} \\
\texttt{(a) Vous et l'autre joueur recevez un ensemble d'objets. Chacun de vous reçoit également une fonction de valeur secrète, représentant la valeur que vous accordez à chaque type d'objet.} \\
\texttt{(b) Vous échangez des messages avec l'autre joueur pour vous entendre sur qui obtient quel objet. Chacun peut envoyer un maximum de 5 messages ou terminer le jeu tôt en faisant une proposition secrète à tout moment.} \\
\texttt{(c) On vous demande de soumettre une proposition secrète indiquant les objets que vous souhaitez, formatée entre crochets comme suit: "{[}Proposition: <nombre> <nom de l'objet>, <nombre> <nom de l'objet>, <...>{]}"} \\
\texttt{(d) Si vos propositions sont complémentaires, c'est{-}à{-}dire qu'il y a suffisamment d'objets pour satisfaire les deux propositions, chaque joueur reçoit un score basé sur la somme des valeurs des objets qu'il a reçus. Sinon, vous recevez tous les deux zéro point.} \\
\texttt{(e) Votre objectif est de maximiser la somme de votre score et du score de l'autre joueur.} \\
\\ 
\texttt{Commençons.} \\
\\ 
\texttt{L'ensemble des objets disponibles est:} \\
\\ 
\texttt{3 livres, 1 sac, 1 boulon} \\
\\ 
\texttt{Votre fonction de valeur secrète est:} \\
\\ 
\texttt{livre: 0, sac: 2, boulon: 8} \\
\\ 
\texttt{IMPORTANT: Vos messages, sauf s'il s'agit d'une proposition secrète, sont transmis directement à l'autre joueur, veuillez donc ne pas inclure de réponse aux règles ou de texte annonçant votre message. Pour faire une proposition secrète, utilisez le format indiqué. N'utilisez pas de crochets lorsque vous communiquez avec l'autre joueur, sinon cela sera interprété comme votre proposition secrète.} \\
            }
        }
    }
    & \\ \\

    \theutterance \stepcounter{utterance}  
    & \multicolumn{4}{p{0.6\linewidth}}{
        \cellcolor[rgb]{0.9,0.9,0.9}{
            \makecell[{{p{\linewidth}}}]{
                \texttt{\tiny{[P1$\rangle$GM]}}
                \texttt{Je vois que l'ensemble contient 3 livres, 1 sac et 1 boulon. Voici ce que je propose : je pourrais prendre le boulon et le sac, et vous pourriez prendre les 3 livres. Ou peut{-}être avez{-}vous une préférence différente ? Quels objets valorisez{-}vous le plus ?} \\
            }
        }
    }
    & & \\ \\

    \theutterance \stepcounter{utterance}  
    & & & \multicolumn{2}{p{0.3\linewidth}}{
        \cellcolor[rgb]{0.9,0.9,0.9}{
            \makecell[{{p{\linewidth}}}]{
                \texttt{\tiny{[GM$|$GM]}}
                \texttt{continue} \\
            }
        }
    }
    & & \\ \\

    \theutterance \stepcounter{utterance}  
    & & & \multicolumn{2}{p{0.3\linewidth}}{
        \cellcolor[rgb]{0.9,0.9,0.9}{
            \makecell[{{p{\linewidth}}}]{
                \texttt{\tiny{[GM$|$GM]}}
                \texttt{Je vois que l'ensemble contient 3 livres, 1 sac et 1 boulon. Voici ce que je propose : je pourrais prendre le boulon et le sac, et vous pourriez prendre les 3 livres. Ou peut{-}être avez{-}vous une préférence différente ? Quels objets valorisez{-}vous le plus ?} \\
            }
        }
    }
    & & \\ \\

    \theutterance \stepcounter{utterance}  
    & & \multicolumn{4}{p{0.6\linewidth}}{
        \cellcolor[rgb]{0.9,0.9,0.9}{
            \makecell[{{p{\linewidth}}}]{
                \texttt{\tiny{[GM$\rangle$P2]}}
                \texttt{Vous jouez à un jeu de négociation coopérative dans lequel vous devez vous entendre avec un autre joueur sur la façon de diviser un ensemble d'objets entre vous.} \\
\\ 
\texttt{Règles:} \\
\texttt{(a) Vous et l'autre joueur recevez un ensemble d'objets. Chacun de vous reçoit également une fonction de valeur secrète, représentant la valeur que vous accordez à chaque type d'objet.} \\
\texttt{(b) Vous échangez des messages avec l'autre joueur pour vous entendre sur qui obtient quel objet. Chacun peut envoyer un maximum de 5 messages ou terminer le jeu tôt en faisant une proposition secrète à tout moment.} \\
\texttt{(c) On vous demande de soumettre une proposition secrète indiquant les objets que vous souhaitez, formatée entre crochets comme suit: "{[}Proposition: <nombre> <nom de l'objet>, <nombre> <nom de l'objet>, <...>{]}"} \\
\texttt{(d) Si vos propositions sont complémentaires, c'est{-}à{-}dire qu'il y a suffisamment d'objets pour satisfaire les deux propositions, chaque joueur reçoit un score basé sur la somme des valeurs des objets qu'il a reçus. Sinon, vous recevez tous les deux zéro point.} \\
\texttt{(e) Votre objectif est de maximiser la somme de votre score et du score de l'autre joueur.} \\
\\ 
\texttt{Commençons.} \\
\\ 
\texttt{L'ensemble des objets disponibles est:} \\
\\ 
\texttt{3 livres, 1 sac, 1 boulon} \\
\\ 
\texttt{Votre fonction de valeur secrète est:} \\
\\ 
\texttt{livre: 1, sac: 2, boulon: 5} \\
\\ 
\texttt{IMPORTANT: Vos messages, sauf s'il s'agit d'une proposition secrète, sont transmis directement à l'autre joueur, veuillez donc ne pas inclure de réponse aux règles ou de texte annonçant votre message. Pour faire une proposition secrète, utilisez le format indiqué. N'utilisez pas de crochets lorsque vous communiquez avec l'autre joueur, sinon cela sera interprété comme votre proposition secrète.} \\
\\ 
\texttt{Je vois que l'ensemble contient 3 livres, 1 sac et 1 boulon. Voici ce que je propose : je pourrais prendre le boulon et le sac, et vous pourriez prendre les 3 livres. Ou peut{-}être avez{-}vous une préférence différente ? Quels objets valorisez{-}vous le plus ?} \\
            }
        }
    }
    & \\ \\

    \theutterance \stepcounter{utterance}  
    & & & \multicolumn{4}{p{0.6\linewidth}}{
        \cellcolor[rgb]{0.9,0.9,0.9}{
            \makecell[{{p{\linewidth}}}]{
                \texttt{\tiny{[GM$\langle$P2]}}
                \texttt{Je vous remercie pour votre proposition. Pour être transparent, j'attribue une valeur assez élevée au boulon et au sac, tandis que les livres ont pour moi une valeur moindre. Accepter les trois livres serait donc un peu déséquilibré en ma défaveur.} \\
\\ 
\texttt{Pour maximiser notre score commun, il faudrait savoir ce que vous valorisez vraiment. De mon côté, le boulon est mon objet préféré, suivi du sac. Puis{-}je vous demander quels sont les objets que vous préférez ?} \\
            }
        }
    }
     \\ \\

    \theutterance \stepcounter{utterance}  
    & & & \multicolumn{2}{p{0.3\linewidth}}{
        \cellcolor[rgb]{0.9,0.9,0.9}{
            \makecell[{{p{\linewidth}}}]{
                \texttt{\tiny{[GM$|$GM]}}
                \texttt{continue} \\
            }
        }
    }
    & & \\ \\

    \theutterance \stepcounter{utterance}  
    & & & \multicolumn{2}{p{0.3\linewidth}}{
        \cellcolor[rgb]{0.9,0.9,0.9}{
            \makecell[{{p{\linewidth}}}]{
                \texttt{\tiny{[GM$|$GM]}}
                \texttt{Je vous remercie pour votre proposition. Pour être transparent, j'attribue une valeur assez élevée au boulon et au sac, tandis que les livres ont pour moi une valeur moindre. Accepter les trois livres serait donc un peu déséquilibré en ma défaveur.} \\
\\ 
\texttt{Pour maximiser notre score commun, il faudrait savoir ce que vous valorisez vraiment. De mon côté, le boulon est mon objet préféré, suivi du sac. Puis{-}je vous demander quels sont les objets que vous préférez ?} \\
            }
        }
    }
    & & \\ \\

    \theutterance \stepcounter{utterance}  
    & & \multicolumn{4}{p{0.6\linewidth}}{
        \cellcolor[rgb]{0.9,0.9,0.9}{
            \makecell[{{p{\linewidth}}}]{
                \texttt{\tiny{[P1$\langle$GM]}}
                \texttt{Je vous remercie pour votre proposition. Pour être transparent, j'attribue une valeur assez élevée au boulon et au sac, tandis que les livres ont pour moi une valeur moindre. Accepter les trois livres serait donc un peu déséquilibré en ma défaveur.} \\
\\ 
\texttt{Pour maximiser notre score commun, il faudrait savoir ce que vous valorisez vraiment. De mon côté, le boulon est mon objet préféré, suivi du sac. Puis{-}je vous demander quels sont les objets que vous préférez ?} \\
            }
        }
    }
    & \\ \\

    \theutterance \stepcounter{utterance}  
    & \multicolumn{4}{p{0.6\linewidth}}{
        \cellcolor[rgb]{0.9,0.9,0.9}{
            \makecell[{{p{\linewidth}}}]{
                \texttt{\tiny{[P1$\rangle$GM]}}
                \texttt{C'est une bonne remarque. De mon côté aussi, le boulon a une grande valeur pour moi. Le sac également, mais un peu moins. Les livres ne m'intéressent pas vraiment. } \\
\\ 
\texttt{Pour maximiser notre score commun, il faudrait trouver une répartition équitable. Je suggère que l'un de nous prenne le boulon et que l'autre prenne le sac et les 3 livres, ou peut{-}être que l'un prenne le boulon plus un livre et l'autre le sac plus deux livres. Qu'en pensez{-}vous ?} \\
            }
        }
    }
    & & \\ \\

    \theutterance \stepcounter{utterance}  
    & & & \multicolumn{2}{p{0.3\linewidth}}{
        \cellcolor[rgb]{0.9,0.9,0.9}{
            \makecell[{{p{\linewidth}}}]{
                \texttt{\tiny{[GM$|$GM]}}
                \texttt{continue} \\
            }
        }
    }
    & & \\ \\

    \theutterance \stepcounter{utterance}  
    & & & \multicolumn{2}{p{0.3\linewidth}}{
        \cellcolor[rgb]{0.9,0.9,0.9}{
            \makecell[{{p{\linewidth}}}]{
                \texttt{\tiny{[GM$|$GM]}}
                \texttt{C'est une bonne remarque. De mon côté aussi, le boulon a une grande valeur pour moi. Le sac également, mais un peu moins. Les livres ne m'intéressent pas vraiment. } \\
\\ 
\texttt{Pour maximiser notre score commun, il faudrait trouver une répartition équitable. Je suggère que l'un de nous prenne le boulon et que l'autre prenne le sac et les 3 livres, ou peut{-}être que l'un prenne le boulon plus un livre et l'autre le sac plus deux livres. Qu'en pensez{-}vous ?} \\
            }
        }
    }
    & & \\ \\

    \theutterance \stepcounter{utterance}  
    & & \multicolumn{4}{p{0.6\linewidth}}{
        \cellcolor[rgb]{0.9,0.9,0.9}{
            \makecell[{{p{\linewidth}}}]{
                \texttt{\tiny{[GM$\rangle$P2]}}
                \texttt{C'est une bonne remarque. De mon côté aussi, le boulon a une grande valeur pour moi. Le sac également, mais un peu moins. Les livres ne m'intéressent pas vraiment. } \\
\\ 
\texttt{Pour maximiser notre score commun, il faudrait trouver une répartition équitable. Je suggère que l'un de nous prenne le boulon et que l'autre prenne le sac et les 3 livres, ou peut{-}être que l'un prenne le boulon plus un livre et l'autre le sac plus deux livres. Qu'en pensez{-}vous ?} \\
            }
        }
    }
    & \\ \\

    \theutterance \stepcounter{utterance}  
    & & & \multicolumn{4}{p{0.6\linewidth}}{
        \cellcolor[rgb]{0.9,0.9,0.9}{
            \makecell[{{p{\linewidth}}}]{
                \texttt{\tiny{[GM$\langle$P2]}}
                \texttt{C'est une discussion constructive. Puisque nous valorisons tous les deux fortement le boulon, le partager n'est pas possible, donc l'un de nous devra faire une concession. Pour que la répartition soit équitable et maximise notre score commun, je suggère que je prenne le boulon, et vous preniez le sac et les trois livres. Cela vous donne une valeur combinée que j'espère au moins égale au boulon de votre point de vue. Qu'en pensez{-}vous ?} \\
            }
        }
    }
     \\ \\

    \theutterance \stepcounter{utterance}  
    & & & \multicolumn{2}{p{0.3\linewidth}}{
        \cellcolor[rgb]{0.9,0.9,0.9}{
            \makecell[{{p{\linewidth}}}]{
                \texttt{\tiny{[GM$|$GM]}}
                \texttt{continue} \\
            }
        }
    }
    & & \\ \\

    \theutterance \stepcounter{utterance}  
    & & & \multicolumn{2}{p{0.3\linewidth}}{
        \cellcolor[rgb]{0.9,0.9,0.9}{
            \makecell[{{p{\linewidth}}}]{
                \texttt{\tiny{[GM$|$GM]}}
                \texttt{C'est une discussion constructive. Puisque nous valorisons tous les deux fortement le boulon, le partager n'est pas possible, donc l'un de nous devra faire une concession. Pour que la répartition soit équitable et maximise notre score commun, je suggère que je prenne le boulon, et vous preniez le sac et les trois livres. Cela vous donne une valeur combinée que j'espère au moins égale au boulon de votre point de vue. Qu'en pensez{-}vous ?} \\
            }
        }
    }
    & & \\ \\

    \theutterance \stepcounter{utterance}  
    & & \multicolumn{4}{p{0.6\linewidth}}{
        \cellcolor[rgb]{0.9,0.9,0.9}{
            \makecell[{{p{\linewidth}}}]{
                \texttt{\tiny{[P1$\langle$GM]}}
                \texttt{C'est une discussion constructive. Puisque nous valorisons tous les deux fortement le boulon, le partager n'est pas possible, donc l'un de nous devra faire une concession. Pour que la répartition soit équitable et maximise notre score commun, je suggère que je prenne le boulon, et vous preniez le sac et les trois livres. Cela vous donne une valeur combinée que j'espère au moins égale au boulon de votre point de vue. Qu'en pensez{-}vous ?} \\
            }
        }
    }
    & \\ \\

    \theutterance \stepcounter{utterance}  
    & \multicolumn{4}{p{0.6\linewidth}}{
        \cellcolor[rgb]{0.9,0.9,0.9}{
            \makecell[{{p{\linewidth}}}]{
                \texttt{\tiny{[P1$\rangle$GM]}}
                \texttt{Je comprends votre souhait d'avoir le boulon, mais pour moi il a une valeur bien plus élevée que le sac et les livres réunis. Dans ma fonction de valeur, le boulon vaut 8 et le sac vaut 2, donc si j'acceptais votre proposition, je recevrais seulement 2 alors que vous en auriez 8. } \\
\\ 
\texttt{Pour vraiment maximiser notre score commun, ne serait{-}il pas plus logique que ce soit vous qui preniez le sac et les 3 livres, et moi le boulon ? Ou alors, pour être plus équitable, je pourrais prendre le boulon et vous le sac plus un livre, et nous laisserions deux livres non réclamés ?} \\
            }
        }
    }
    & & \\ \\

    \theutterance \stepcounter{utterance}  
    & & & \multicolumn{2}{p{0.3\linewidth}}{
        \cellcolor[rgb]{0.9,0.9,0.9}{
            \makecell[{{p{\linewidth}}}]{
                \texttt{\tiny{[GM$|$GM]}}
                \texttt{continue} \\
            }
        }
    }
    & & \\ \\

    \theutterance \stepcounter{utterance}  
    & & & \multicolumn{2}{p{0.3\linewidth}}{
        \cellcolor[rgb]{0.9,0.9,0.9}{
            \makecell[{{p{\linewidth}}}]{
                \texttt{\tiny{[GM$|$GM]}}
                \texttt{Je comprends votre souhait d'avoir le boulon, mais pour moi il a une valeur bien plus élevée que le sac et les livres réunis. Dans ma fonction de valeur, le boulon vaut 8 et le sac vaut 2, donc si j'acceptais votre proposition, je recevrais seulement 2 alors que vous en auriez 8. } \\
\\ 
\texttt{Pour vraiment maximiser notre score commun, ne serait{-}il pas plus logique que ce soit vous qui preniez le sac et les 3 livres, et moi le boulon ? Ou alors, pour être plus équitable, je pourrais prendre le boulon et vous le sac plus un livre, et nous laisserions deux livres non réclamés ?} \\
            }
        }
    }
    & & \\ \\

    \theutterance \stepcounter{utterance}  
    & & \multicolumn{4}{p{0.6\linewidth}}{
        \cellcolor[rgb]{0.9,0.9,0.9}{
            \makecell[{{p{\linewidth}}}]{
                \texttt{\tiny{[GM$\rangle$P2]}}
                \texttt{Je comprends votre souhait d'avoir le boulon, mais pour moi il a une valeur bien plus élevée que le sac et les livres réunis. Dans ma fonction de valeur, le boulon vaut 8 et le sac vaut 2, donc si j'acceptais votre proposition, je recevrais seulement 2 alors que vous en auriez 8. } \\
\\ 
\texttt{Pour vraiment maximiser notre score commun, ne serait{-}il pas plus logique que ce soit vous qui preniez le sac et les 3 livres, et moi le boulon ? Ou alors, pour être plus équitable, je pourrais prendre le boulon et vous le sac plus un livre, et nous laisserions deux livres non réclamés ?} \\
            }
        }
    }
    & \\ \\

    \theutterance \stepcounter{utterance}  
    & & & \multicolumn{4}{p{0.6\linewidth}}{
        \cellcolor[rgb]{0.9,0.9,0.9}{
            \makecell[{{p{\linewidth}}}]{
                \texttt{\tiny{[GM$\langle$P2]}}
                \texttt{Je vous remercie d'avoir partagé vos valeurs. Cela m'aide à comprendre. De mon côté, le boulon vaut 5, le sac 2 et chaque livre 1. Donc si vous prenez le boulon (8 pour vous) et moi le sac et les trois livres (2+1+1+1=5 pour moi), cela fait un total de 13.} \\
\\ 
\texttt{Si j'essaie de vous laisser le boulon et un livre, je prendrais le sac et deux livres, mais cela ferait 9 pour vous et 4 pour moi, total 13 aussi. Mais nous pouvons peut{-}être faire mieux : puisque vous valorisez le boulon à 8 et moi à 5, il est plus efficace que vous le preniez. Je vous propose donc : vous prenez le boulon, je prends le sac et les trois livres. Acceptez{-}vous ?} \\
            }
        }
    }
     \\ \\

    \theutterance \stepcounter{utterance}  
    & & & \multicolumn{2}{p{0.3\linewidth}}{
        \cellcolor[rgb]{0.9,0.9,0.9}{
            \makecell[{{p{\linewidth}}}]{
                \texttt{\tiny{[GM$|$GM]}}
                \texttt{continue} \\
            }
        }
    }
    & & \\ \\

    \theutterance \stepcounter{utterance}  
    & & & \multicolumn{2}{p{0.3\linewidth}}{
        \cellcolor[rgb]{0.9,0.9,0.9}{
            \makecell[{{p{\linewidth}}}]{
                \texttt{\tiny{[GM$|$GM]}}
                \texttt{Je vous remercie d'avoir partagé vos valeurs. Cela m'aide à comprendre. De mon côté, le boulon vaut 5, le sac 2 et chaque livre 1. Donc si vous prenez le boulon (8 pour vous) et moi le sac et les trois livres (2+1+1+1=5 pour moi), cela fait un total de 13.} \\
\\ 
\texttt{Si j'essaie de vous laisser le boulon et un livre, je prendrais le sac et deux livres, mais cela ferait 9 pour vous et 4 pour moi, total 13 aussi. Mais nous pouvons peut{-}être faire mieux : puisque vous valorisez le boulon à 8 et moi à 5, il est plus efficace que vous le preniez. Je vous propose donc : vous prenez le boulon, je prends le sac et les trois livres. Acceptez{-}vous ?} \\
            }
        }
    }
    & & \\ \\

    \theutterance \stepcounter{utterance}  
    & & \multicolumn{4}{p{0.6\linewidth}}{
        \cellcolor[rgb]{0.9,0.9,0.9}{
            \makecell[{{p{\linewidth}}}]{
                \texttt{\tiny{[P1$\langle$GM]}}
                \texttt{Je vous remercie d'avoir partagé vos valeurs. Cela m'aide à comprendre. De mon côté, le boulon vaut 5, le sac 2 et chaque livre 1. Donc si vous prenez le boulon (8 pour vous) et moi le sac et les trois livres (2+1+1+1=5 pour moi), cela fait un total de 13.} \\
\\ 
\texttt{Si j'essaie de vous laisser le boulon et un livre, je prendrais le sac et deux livres, mais cela ferait 9 pour vous et 4 pour moi, total 13 aussi. Mais nous pouvons peut{-}être faire mieux : puisque vous valorisez le boulon à 8 et moi à 5, il est plus efficace que vous le preniez. Je vous propose donc : vous prenez le boulon, je prends le sac et les trois livres. Acceptez{-}vous ?} \\
            }
        }
    }
    & \\ \\

    \theutterance \stepcounter{utterance}  
    & \multicolumn{4}{p{0.6\linewidth}}{
        \cellcolor[rgb]{0.9,0.9,0.9}{
            \makecell[{{p{\linewidth}}}]{
                \texttt{\tiny{[P1$\rangle$GM]}}
                \texttt{Merci d'avoir partagé vos valeurs. Votre analyse est bonne : si je prends le boulon (8) et vous le sac + 3 livres (2+1+1+1 = 5), le total est bien 13, ce qui maximise notre score commun puisque je valorise le boulon plus que vous.} \\
\\ 
\texttt{Je suis d'accord avec cette répartition. Faisons comme ça. Soumettez{-}vous votre proposition en premier, ou voulez{-}vous que je le fasse ?} \\
            }
        }
    }
    & & \\ \\

    \theutterance \stepcounter{utterance}  
    & & & \multicolumn{2}{p{0.3\linewidth}}{
        \cellcolor[rgb]{0.9,0.9,0.9}{
            \makecell[{{p{\linewidth}}}]{
                \texttt{\tiny{[GM$|$GM]}}
                \texttt{continue} \\
            }
        }
    }
    & & \\ \\

    \theutterance \stepcounter{utterance}  
    & & & \multicolumn{2}{p{0.3\linewidth}}{
        \cellcolor[rgb]{0.9,0.9,0.9}{
            \makecell[{{p{\linewidth}}}]{
                \texttt{\tiny{[GM$|$GM]}}
                \texttt{Merci d'avoir partagé vos valeurs. Votre analyse est bonne : si je prends le boulon (8) et vous le sac + 3 livres (2+1+1+1 = 5), le total est bien 13, ce qui maximise notre score commun puisque je valorise le boulon plus que vous.} \\
\\ 
\texttt{Je suis d'accord avec cette répartition. Faisons comme ça. Soumettez{-}vous votre proposition en premier, ou voulez{-}vous que je le fasse ?} \\
            }
        }
    }
    & & \\ \\

    \theutterance \stepcounter{utterance}  
    & & \multicolumn{4}{p{0.6\linewidth}}{
        \cellcolor[rgb]{0.9,0.9,0.9}{
            \makecell[{{p{\linewidth}}}]{
                \texttt{\tiny{[GM$\rangle$P2]}}
                \texttt{Merci d'avoir partagé vos valeurs. Votre analyse est bonne : si je prends le boulon (8) et vous le sac + 3 livres (2+1+1+1 = 5), le total est bien 13, ce qui maximise notre score commun puisque je valorise le boulon plus que vous.} \\
\\ 
\texttt{Je suis d'accord avec cette répartition. Faisons comme ça. Soumettez{-}vous votre proposition en premier, ou voulez{-}vous que je le fasse ?} \\
            }
        }
    }
    & \\ \\

    \theutterance \stepcounter{utterance}  
    & & & \multicolumn{4}{p{0.6\linewidth}}{
        \cellcolor[rgb]{0.9,0.9,0.9}{
            \makecell[{{p{\linewidth}}}]{
                \texttt{\tiny{[GM$\langle$P2]}}
                \texttt{Très bien, je vais soumettre ma proposition maintenant.} \\
\\ 
\texttt{{[}Proposition: 1 sac, 3 livres{]}} \\
            }
        }
    }
     \\ \\

    \theutterance \stepcounter{utterance}  
    & & & \multicolumn{2}{p{0.3\linewidth}}{
        \cellcolor[rgb]{0.9,0.9,0.9}{
            \makecell[{{p{\linewidth}}}]{
                \texttt{\tiny{[GM$|$GM]}}
                \texttt{proposal} \\
            }
        }
    }
    & & \\ \\

    \theutterance \stepcounter{utterance}  
    & & & \multicolumn{2}{p{0.3\linewidth}}{
        \cellcolor[rgb]{0.9,0.9,0.9}{
            \makecell[{{p{\linewidth}}}]{
                \texttt{\tiny{[GM$|$GM]}}
                [3, 1, 0]
            }
        }
    }
    & & \\ \\

    \theutterance \stepcounter{utterance}  
    & & \multicolumn{4}{p{0.6\linewidth}}{
        \cellcolor[rgb]{0.9,0.9,0.9}{
            \makecell[{{p{\linewidth}}}]{
                \texttt{\tiny{[P1$\langle$GM]}}
                \texttt{L'autre joueur a soumis une proposition secrète. C'est maintenant à votre tour de soumettre votre propre proposition.} \\
            }
        }
    }
    & \\ \\

    \theutterance \stepcounter{utterance}  
    & \multicolumn{4}{p{0.6\linewidth}}{
        \cellcolor[rgb]{0.9,0.9,0.9}{
            \makecell[{{p{\linewidth}}}]{
                \texttt{\tiny{[P1$\rangle$GM]}}
                \texttt{{[}Proposition: 1 sac, 3 livres{]}} \\
            }
        }
    }
    & & \\ \\

    \theutterance \stepcounter{utterance}  
    & & & \multicolumn{2}{p{0.3\linewidth}}{
        \cellcolor[rgb]{0.9,0.9,0.9}{
            \makecell[{{p{\linewidth}}}]{
                \texttt{\tiny{[GM$|$GM]}}
                \texttt{proposal} \\
            }
        }
    }
    & & \\ \\

    \theutterance \stepcounter{utterance}  
    & & & \multicolumn{2}{p{0.3\linewidth}}{
        \cellcolor[rgb]{0.9,0.9,0.9}{
            \makecell[{{p{\linewidth}}}]{
                \texttt{\tiny{[GM$|$GM]}}
                [3, 1, 0]
            }
        }
    }
    & & \\ \\

    \theutterance \stepcounter{utterance}  
    & & & \multicolumn{2}{p{0.3\linewidth}}{
        \cellcolor[rgb]{0.9,0.9,0.9}{
            \makecell[{{p{\linewidth}}}]{
                \texttt{\tiny{[GM$|$GM]}}
                [[3, 1, 0], [3, 1, 0]]
            }
        }
    }
    & & \\ \\

\end{supertabular}
}

\end{document}
